\usetikzlibrary{arrows,decorations.pathmorphing,backgrounds,positioning,fit,shapes.misc,chains,scopes,shadows}

\tikzset{
  syntax-rule/.style={
    % The shape: 
    rectangle, 
    % The size: 
    minimum size=5mm,
    minimum width=8em, 
    % The border: 
    % very thick, 
    % draw=red!50!black!50,       % 50% red and 50% black, 
                                % and that mixed with 50% white
    % The filling: 
    % top color=white, 
    % bottom color=red!50!black!20, % and something else at the bottom 
    % Font 
    font=\itshape\small
  }
}

\tikzset{
  syntax-rule-cont/.style={
    % The shape: 
    rectangle, 
    % The size: 
    minimum size=6mm, 
    minimum width=8em,
    % The border: 
    % very thick, 
    % draw=red!50!black!50,       % 50% red and 50% black, 
                                % and that mixed with 50% white
    % The filling: 
    % top color=white, 
    % bottom color=red!50!black!20, % and something else at the bottom 
    % Font 
    font=\itshape\small,
    text=white
  }
}

\tikzset{
  nonterminal/.style={
      % The shape: 
      rectangle, 
      % The size: 
      minimum size=5mm, 
      % The border: 
      very thick, 
      draw=red!30!black!50,       % 50% red and 50% black, 
                                  % and that mixed with 50% white
      % The filling: 
      top color=white, 
      bottom color=red!30!black!20, % and something else at the bottom 
      % Font 
      font=\small
    }
}

\tikzset{
  terminal/.style={
    % The shape:
    rounded rectangle,
    % The size:
    minimum size=5mm,
    % The border:
    very thick,
    draw=black!50,
    % The filling:
    top color=white,
    bottom color=black!20,
    font=\ttfamily\small
  }
}

\tikzset{
  meta-terminal/.style={
    % The shape:
    rounded rectangle,
    % The size:
    minimum size=5mm,
    % The border:
    very thick,
    draw=black!50,
    opacity=0.5,
    % The filling:
    top color=white,
    bottom color=black!20,
    font=\ttfamily\it\small
  }
}


% Define box and box title style


\definecolor{mynote-color}{rgb}{1, 0.8, 0.4}

\tikzstyle{mynote-mybox} = [draw=mynote-color, fill=mynote-color, very thick, fill opacity=0.1,
    rectangle, rounded corners, inner sep=10pt]
\tikzstyle{mynote-fancytitle} =[fill=mynote-color, text=white]


% =============
% = C Section =
% =============

\definecolor{c-color}{rgb}{0.5, 0.05, 0}

\tikzstyle{c-mybox} = [draw=c-color, fill=c-color, very thick, fill opacity=0.05, rectangle, rounded corners, inner sep=10pt]
\tikzstyle{c-fancytitle} =[fill=c-color, text=white]

\newcommand{\csection}[1] {%
\begin{tikzpicture}%
\node [c-mybox] (box){%
    \begin{minipage}{0.95\textwidth}%
      \pgfsetfillopacity{1}%
      #1%
    \end{minipage}%
};%
\node[c-fancytitle, right=10pt] at (box.north west) { C };
\node[c-fancytitle, rounded corners] at (box.east) {$\diamondsuit$};
\end{tikzpicture}%
}

\newcommand{\lcsection}[1] {%
\clearpage
\thispagestyle{codepage}
\begin{tikzpicture}[remember picture, overlay]%
\node [c-mybox,minimum width=0.90\paperheight,minimum height=0.90\paperwidth,rotate=90] (box) at (current page.center) {%
    \begin{minipage}{\textheight}%
      \pgfsetfillopacity{1}%
      #1%
    \end{minipage}%
};%
\node[c-fancytitle, rotate=90] at (box.north) {C};
\node[c-fancytitle, rotate=90, rounded corners] at (box.east) {$\diamondsuit$};
\end{tikzpicture}%
\clearpage
}


% ==================
% = Pascal Section =
% ==================

\tikzstyle{pas-mybox} = [draw=blue!50, fill=blue!20, very thick, fill opacity=0.1, rectangle, rounded corners, inner sep=10pt]
\tikzstyle{pas-fancytitle} =[fill=blue!50, text=white]


\newcommand{\passection}[1] {%
\begin{tikzpicture}%
\node [pas-mybox] (box){%
    \begin{minipage}{0.95\textwidth}%
      \pgfsetfillopacity{1}%
      #1%
    \end{minipage}%
};%
\node[pas-fancytitle, right=10pt] at (box.north west) {Pascal};
\node[pas-fancytitle, rounded corners] at (box.east) {$\heartsuit$};
\end{tikzpicture}%
}

\newcommand{\lpassection}[1] {%
\clearpage
\thispagestyle{codepage}
\begin{tikzpicture}[remember picture, overlay]%
\node [pas-mybox,minimum width=0.90\paperheight,minimum height=0.90\paperwidth,rotate=90] (box) at (current page.center) {%
    \begin{minipage}{\textheight}%
      \pgfsetfillopacity{1}%
      #1%
    \end{minipage}%
};%
\node[pas-fancytitle, rotate=90] at (box.north) {Pascal};
\node[pas-fancytitle, rotate=90, rounded corners] at (box.east) {$\heartsuit$};
\end{tikzpicture}%
\clearpage
}



% ===========
% = My Note =
% ===========

\newcommand{\mynote}[1] {%
\begin{tikzpicture}%
\node [mynote-mybox] (box){%
    \begin{minipage}{0.95\textwidth}%
      \pgfsetfillopacity{1}%  
      #1%
    \end{minipage}%
};%
\node[mynote-fancytitle, right=10pt] at (box.north west) {Note};
\node[mynote-fancytitle, rounded corners] at (box.east) {$\spadesuit$};
\end{tikzpicture}%
}


% =============================================
% = Pseudocode - pseudocode and pseudosection =
% =============================================
\definecolor{pseudo-color}{rgb}{0.8, 0.4, 1}

\tikzstyle{pseudo-mybox} = [draw=pseudo-color, fill=pseudo-color, very thick, fill opacity=0.2,
    rectangle, rounded corners, inner sep=10pt]
\tikzstyle{pseudo-fancytitle} =[fill=pseudo-color, text=white]

\newcommand{\pseudosection}[1] {%
\begin{tikzpicture}%
\node [pseudo-mybox] (box){%
    \begin{minipage}{0.95\textwidth}%
      \pgfsetfillopacity{1}%      
        #1 %
    \end{minipage}%
};%
\node[pseudo-fancytitle, right=10pt] at (box.north west) {Pseudocode};
\node[pseudo-fancytitle, rounded corners] at (box.east) {$\clubsuit$};
\end{tikzpicture}%
}

\newcommand{\lpseudosection}[1] {%
\clearpage
\thispagestyle{codepage}
\begin{tikzpicture}[remember picture, overlay]%
\node [pseudo-mybox,minimum width=0.90\paperheight,minimum height=0.90\paperwidth,rotate=90] (box) at (current page.center) {%
    \begin{minipage}{\textheight}%
      \pgfsetfillopacity{1}%
      #1%
    \end{minipage}%
};%
\node[pseudo-fancytitle, rotate=90] at (box.north) {Pseudocode};
\node[pseudo-fancytitle, rotate=90, rounded corners] at (box.east) {$\clubsuit$};
\end{tikzpicture}%
\clearpage
}

\newcommand{\pseudocode}[3]
{
  \lstset{language=Pascal}
  \pseudosection{\lstinputlisting[caption={#2},label={#1}]{#3}}
}


\newcommand{\lpseudocode}[4]
{
  \lstset{language=Pascal}
  \lpseudosection{\lstinputlisting[caption={#2},label={#1},multicols=#4]{#3}}
}

% =======================
% = Syntax - syntaxsect =
% =======================

\newcommand{\csyntax}[3] {%
  \begin{figure}[htbp]%
   \centering%
    
    \begin{tikzpicture}
      \draw[very thick,color=gray] (0,-0.15) -- (0,0) -- (\textwidth,0) -- (\textwidth,-0.15);
    \end{tikzpicture}
    
    %\frame{
    \begin{minipage}{0.95\textwidth}%
      \input{./syntax-out/c/#3}%
      \smallskip
    \end{minipage}%
    %}
    
    \begin{tikzpicture}
      \draw[very thick,color=gray] (0,0.15) -- (0,0) -- (\textwidth,0) -- (\textwidth,0.15);
    \end{tikzpicture}
    \caption{C Syntax for #2}%
    \label{#1}%
  \end{figure}%
}

\newcommand{\passyntax}[3] {%
  \begin{figure}[htbp]%
   \centering%
    
    \begin{tikzpicture}
      \draw[very thick,color=gray] (0,-0.15) -- (0,0) -- (\textwidth,0) -- (\textwidth,-0.15);
    \end{tikzpicture}
    
    %\frame{
    \begin{minipage}{0.95\textwidth}%
      \input{./syntax-out/pascal/#3}%
      \smallskip
    \end{minipage}%
    %}
    
    \begin{tikzpicture}
      \draw[very thick,color=gray] (0,0.15) -- (0,0) -- (\textwidth,0) -- (\textwidth,0.15);
    \end{tikzpicture}
    \caption{Pascal Syntax for #2}%
    \label{#1}%
  \end{figure}%
}

\newcommand{\examplesyntax}[3] {%
  \begin{figure}[htbp]%
   \centering%
    
    \begin{tikzpicture}
      \draw[very thick,color=gray] (0,-0.15) -- (0,0) -- (\textwidth,0) -- (\textwidth,-0.15);
    \end{tikzpicture}
    
    %\frame{
    \begin{minipage}{0.95\textwidth}%
      \input{./syntax-out/example/#3}%
      \smallskip
    \end{minipage}%
    %}
    
    \begin{tikzpicture}
      \draw[very thick,color=gray] (0,0.15) -- (0,0) -- (\textwidth,0) -- (\textwidth,0.15);
    \end{tikzpicture}
    \caption{#2}%
    \label{#1}%
  \end{figure}%
}




